

\section*{Supplementary Material}

\section{Marginal Shapley values} 

Marginal Shapley values are heavily influenced by the distribution of the data.

\subsection{Shifted distributions}
\label{sec:suppl_marg_mean}

First, let us consider the linear model $\blackbox(\localobs)=\localobs_1+\localobs_2$ with $\localobs_1\sim\Uniform(-1, 2)$ and $\localobs_2\sim\Uniform(0, 3)$ and local observation $\localobs=(0, 0)$. We show that although they play symmetric roles in the algebraic formulation of the black box model, their marginal values aren't equal: $\phi_{1}=-0.5$ and $\phi_{2}=-1.5$. 

\begin{align*}
    \phi_{1} &= \sum_{S\subseteq \{\emptyset, x_{2}\}} \frac{|S| !(|F|-|S|-1) !}{|F| !} \\
    &\left[f_{S \cup\{i\}}\left(x_{S \cup\{i\}}\right)-f_{S}\left(x_{S}\right)\right]\\
    & = \frac{1}{2} (f_{1}(x_{1})-f_{\emptyset}(\emptyset)) + \frac{1}{2} (f_{1, 2}(x_{1}, x_{2}) - f_{2}(x_{2}))\\
    & \overset{*}{=} \frac{1}{2}E[f(x_{1}, X_{2})] - \frac{1}{2}E[f(X_{1}, X_{2})] \\ &+ \frac{1}{2} E[f(x_{1}, x_{2})] - \frac{1}{2} E[f(X_{1}, x_{2})]\\
    & = \frac{1}{2}P(X_{2}=1) - \frac{1}{2}P(X_{2}=1|X_{1}=1)P(X_{1}=1) \\ &+ \frac{1}{2} - \frac{1}{2} P(X_{1}=1)\\
    & = \frac{1}{2}p - \frac{1}{2}p\cdot 1 + \frac{1}{2} - \frac{1}{2}\cdot 1 = 0
\end{align*}
where we used the definition $f(x_S)=E[f(x)|do\;x_S]$ from \cite{janzing2020feature} for marginal Shapley values in equation $*$.
\begin{align*}
    2\phi_1 & = \frac{1}{3}[\int_{0}^{3} x_2 d x_2 ] - \frac{1}{9} [\int_{0}^{3} \int_{-1}^{2} (x_1+x_2) d x_2 d x_1 ] \\& - \frac{1}{3}[\int_{-1}^{2} x_1 d x_1 ]\\
    & = \frac{1}{3}[\int_{0}^{3} x_2 d x_2 ] - \frac{1}{3} [\int_{-1}^{2} x_1 d d x_1 ] \\& -\frac{1}{3} [\int_{0}^{3} x_2 d x_2 ] - \frac{1}{3}[\int_{-1}^{2} x_1 d x_1 ]\\
    & = - \frac{2}{3} [\int_{-1}^{2} x_1 d x_1 ] \\
    & = - \frac{2}{3} [\int_{-1}^{2} x_1 d x_1 ] \\
    & = - \frac{2}{3} [\frac{4}{2} - \frac{(-1)^2}{2}] \\
\end{align*}
Thus
$$\phi_1 = \frac{-1}{2} $$

Symmetrically for $x_2$,
$$2\phi_2 = - \frac{2}{3} [\int_{0}^{3} x_2 d x_2 ]$$
and therefore
$$\phi_2 = - \frac{3}{2}$$

\subsection{Different spreads}
\label{sec:suppl_marg_spread}

Let us consider a black box $\blackbox(\localobs)=\localobs_1^2+\localobs_2^2$ with $\localobs_1\sim\Normal(0, 1)$ and $\localobs_2\sim\Normal(0, 10)$ and local observation $\localobs=(0, 0)$. While the first marginal Shapley value is -1, the second one is -100 as the expected change of model outcome is higher when intervening on the common population by setting $\localobs_2=0$ compared to setting $\localobs_1=0$. 

Similarly to the previous section:
\begin{align*}
    \phi_1 & = - \int_{-\infty}^{+\infty} \frac{x_{1}^2}{\sqrt{2\pi}} \exp{\frac{-x_{1}^2}2} d x_1 ] \\
    & =\frac{1}{\sqrt{2\pi}}[-x_{1} \mathrm{e}^{-\frac{x_{1}^{2}}{2}}-\int_{-\infty}^{+\infty}-\mathrm{e}^{-\frac{x_{1}^{2}}{2}} \mathrm{~d} x_1 \\
\end{align*}
Solving separately:
$$
\int-\mathrm{e}^{-\frac{x_1^{2}}{2}} \mathrm{~d} x
$$
We substitute $u=\frac{x_1}{\sqrt{2}} \longrightarrow \frac{\mathrm{d} u}{\mathrm{~d} x_1}=\frac{1}{\sqrt{2}}$ $\longrightarrow \mathrm{d} x_1=\sqrt{2} \mathrm{~d} u$ :
$$
=-\frac{\sqrt{\pi}}{\sqrt{2}} \int \frac{2 \mathrm{e}^{-u^{2}}}{\sqrt{\pi}} \mathrm{d} u
$$

We notice the Gaussian error function below: 
$$
\int \frac{2 \mathrm{e}^{-u^{2}}}{\sqrt{\pi}} \mathrm{d} u =\operatorname{erf}(u)
$$
We plug in solved integrals:
$$
\begin{aligned}
&-\frac{\sqrt{\pi}}{\sqrt{2}} \int \frac{2 \mathrm{e}^{-u^{2}}}{\sqrt{\pi}} \mathrm{d} u \\
&=-\frac{\sqrt{\pi} \operatorname{erf}(u)}{\sqrt{2}}
\end{aligned}
$$
We undo the substitution $u=\frac{x}{\sqrt{2}}$ :
$$
=-\frac{\sqrt{\pi} \operatorname{erf}\left(\frac{x}{\sqrt{2}}\right)}{\sqrt{2}}
$$
Ultimately: 
$$
\int_{-\infty}^{+\infty} \frac{x^{2} \mathrm{e}^{-\frac{x^{2}}{2}}}{\sqrt{2} \sqrt{\pi}} \mathrm{d} x =[\frac{\operatorname{erf}\left(\frac{x}{\sqrt{2}}\right)}{2}-\frac{x \mathrm{e}^{-\frac{x^{2}}{2}}}{\sqrt{2} \sqrt{\pi}}]_{-\infty}^{+\infty}=1 
$$\\
Symmetrically for $x_2$,
\begin{align*}
    \phi_2 & = - \int_{-\infty}^{+\infty} \frac{x_{1}^2}{10\sqrt{2\pi}} \exp{\frac{-x_{1}^2}{200}} d x_1 \\
    & = -100\\
\end{align*}

\section{Counterfactual Fairness} \label{sec:suppl_marg_fairness}
Because of the dummy property of marginal Shapley values we have that 
$$ \text{counterfactual fairness} \rightarrow \text{marginal Shapley value of 0}$$
for deterministic models. 
The back direction as we will see does not hold: 
Let our feature space comprise two binary variables $x_{canLift}$ and $x_{male}$ with 
\begin{align*}
    P(x_{male}=1) &= 1 \\
    P(x_{canLift}=1) &= p\\
    P(x_{canLift}=1|x_{male}=1) &= p
\end{align*}
where $p$ is an arbitrary probability.

Our black box algorithm is $f(x)=x_{male}\cdot x_{canLift}$ and the feature attribution of $x_{male}$ for $x_{male}=x_{canLift}=1$ can be computed as follows
\begin{align*}
    & \phi_{male} = \\&  \sum_{S\subseteq \{\emptyset, x_{canLift}\}} \frac{|S| !(|F|-|S|-1) !}{|F| !} [f_{S \cup\{i\}}\left(x_{S \cup\{i\}}\right) \\& -f_{S}\left(x_{S}\right) ]\\
    & = \frac{1}{2} (f_{male}(x_{male})-f_{\emptyset}(\emptyset)) \\& + \frac{1}{2} (f_{male, canLift}(x_{male}, x_{canLift}) - f_{canLift}(x_{canLift}))\\
    & \overset{*}{=} \frac{1}{2}E[f(x_{male}, X_{canLift})] - \frac{1}{2}E[f(X_{male}, X_{canLift})] \\& + \frac{1}{2} E[f(x_{male}, x_{canLift})] - \frac{1}{2} E[f(X_{male}, x_{canLift})]\\
    & = \frac{1}{2}P(X_{canLift}=1) + \frac{1}{2} - \frac{1}{2} P(X_{male}=1) \\& - \frac{1}{2}P(X_{canLift}=1|X_{male}=1)P(X_{male}=1) \\
    & = \frac{1}{2}p - \frac{1}{2}p\cdot 1 + \frac{1}{2} - \frac{1}{2}\cdot 1 = 0
\end{align*}
where we used the definition $f(x_S)=E[f(x)|do\;x_S]$ from \cite{janzing2020feature} for marginal Shapley values in equation $*$.


\section{Feature Selection} 
\label{sec:suppl_featureselect} 
Let our feature space comprise two independent binary variables $\varobs_{1}, \varobs_{2}\sim \Normal(1, 1)$.
Our black box algorithm is defined by $f(x)=\indicate{\localobs_{1}>1}\cdot 3\localobs_{2}-\indicate{\localobs_{1}\leq 1}\cdot \localobs_{2}$ and the conditional feature attribution of feature $2$ at $\localobs=(0.5, 0.5)$ can be computed as follows
\begin{align*}
    \phi_{2} &= \sum_{S\subseteq \{\emptyset, 1\}} \frac{|S| !(|F|-|S|-1) !}{|F| !}[f_{S \cup\{i\}}\left(x_{S \cup\{i\}}\right)
    \\& -f_{S}\left(x_{S}\right)]\\
    & = \frac{1}{2}(f_{2}(x_{2})-f_{\emptyset}(\emptyset)) + \frac{1}{2} (f_{1, 2}(x_{1}, x_{2}) - f_{1}(x_{1}))\\
    & \overset{*}{=} \frac{1}{2}E[f(X_{1}, x_{2}=0.5)] - \frac{1}{2}E[f(X_{1}, X_{2})] \\& + \frac{1}{2} E[f(x_{1}=0.5, x_{2}=0.5)] - \frac{1}{2} E[f(x_{1}=0.5, X_{2})]\\
    & = \frac{1}{2} (0.5\cdot 0.5\cdot 3 -0.5\cdot 0.5 ) - \frac{1}{2}E[0.5\cdot 3X_2-0.5X_2] \\& - \frac{1}{2} 0.5 - \frac{1}{2}E[-X_{2}] \\
    & = \frac{1}{4} - \frac{1}{2} - \frac{1}{4} + \frac{1}{2} = 0
    % & = \frac{1}{2} (0.5\cdot 0.5\cdot 3 -0.5\cdot 0.5 ) - \frac{1}{2}E[0.5\cdot 3X_2-0.5X_2] - \frac{1}{2} 0.5 - \frac{1}{2}E[-X_{2}] \\
    % & = \frac{2}{8} - \frac{1}{4}(a-1) - \frac{1}{4} + \frac{1}{2} \\
\end{align*}
where we used $f_S(x_S)=E[f(x_S,X_{\bar{S}})]$ in equation $*$.

